\section{Threat Model}
\label{sec:threat}

\paragraph{Adversary.}
The adversary $\mathcal{A}$ is a profit-maximizing, computationally bounded party that (i) controls a
budget usable for bribes and gas; (ii) can deploy smart contracts on the sharded system or a
connected chain; (iii) can observe the public mempool, certificates, and committee assignments;
and (iv) may, with bounded probability, induce or exploit transient network asynchrony or occupy a
leader role within the target committee (this is the lever that engineers the honest-vote split
underlying $\psucc$). $\mathcal{A}$ \emph{cannot} break cryptographic primitives, cannot corrupt
more than the network-wide Byzantine threshold across \emph{all} shards simultaneously, and does
not control the randomness used for committee assignment (we treat assignment as given for the
target epoch).

\paragraph{Sub-threshold, precisely.}
$\mathcal{A}$ is \emph{sub-threshold} in the network-wide sense: across the whole validator set it
controls strictly fewer than the global $1/3$ safety threshold. Within the single target committee,
however, it procures $F+1$ equivocating validators---at or above that committee's local threshold.
We never assume $\mathcal{A}$ is below the \emph{local} committee threshold; the entire point is that
the global bound is respected while one committee's bound is not.

\paragraph{Targets and goal.}
$\mathcal{A}$ targets a single committee $C$ of size $N=3F+1$ during one epoch. The goal is to cause
$C$ to finalize two conflicting certificates over shard-local state and to convert the resulting
ambiguity into irreversible value $V$ through the cross-shard or bridge interface, with positive net
profit after bribes, slashing-induced effects, and operating costs. Buying $F+1$ equivocators is
\emph{necessary} (Lemma~\ref{lem:intersection}) but not \emph{sufficient}: $\mathcal{A}$ must also
drive a usable honest-vote split, which succeeds with probability $\psucc$ determined by the protocol
and network (\S\ref{sec:floor}).

\paragraph{Defender.}
The defender runs accountable BFT with a slashing mechanism: provable equivocation by a validator
yields a transferable proof; upon enforcement, the validator's bonded stake $s_i$ is slashed and the
offending branch is reverted or quarantined~\cite{buterin2017casper,buchman2018tendermint}.
Validators have a bonding/unbonding period; let $\Tacc$ denote the end-to-end accountability latency
(detection $\rightarrow$ proof propagation $\rightarrow$ enforcement $\rightarrow$ reversion).
Cross-shard receipts settle with latency $\Tsettle$. We treat a \emph{robust} defender as one whose
stake remains globally bonded and slashable across committee rotation for at least the maximum
evidence delay; \S\ref{sec:reconf} analyzes what happens when this property is weakened.

\paragraph{Assumptions.}
\begin{assumption}[Rational validators]\label{ass:rational}
Each validator maximizes expected payoff and accepts a bribe iff it weakly exceeds the expected cost
of the prescribed behavior.
\end{assumption}
\begin{assumption}[Accountable equivocation]\label{ass:accountable}
Equivocation is detectable and attributable: two validly signed conflicting messages from the same
validator constitute a proof that any honest party can verify and submit.
\end{assumption}
\begin{assumption}[Asynchronous settlement]\label{ass:async}
Cross-shard effects are not gated on the source shard's accountability window unless the protocol
explicitly enforces such gating; absent such a gate, $\Tsettle$ and $\Tacc$ are independent.
\end{assumption}
